\scanpage{16}
\mspage{14}
\uline{pas un chemin} de longueur $m+n$, i.e. un diagramme de type $\Phi_n$, mais un diagramme de type \uline{dual} $\Phi_n^o=\Psi_n$. P.ex. pour $m=n=1$
\diagram{\begin{tikzcd}[ampersand replacement=\&,column sep=.85em,row sep=.2em]
A_0 \& A_1\arrow[l]\arrow[r] \& \underset{\substack{\Vert\\B_0}}{A_2} \& B_1\arrow[l]\arrow[r] \& B_2\\
\& \smash{\underbrace{\hspace{4.8em}}_c} \& \& \smash{\underbrace{\hspace{4.8em}}_{c'}}
\end{tikzcd}}
est de type $\Phi_2^o$ ! Donc pour ne pas trop compliquer la vie, on va dans les développements précédents se borner à $m,n$ \uline{pairs}, donc des chemins de longueur au moins $4$. C'est un peu ridicule ! On est ramené à prouver la commutativité à homotopie près dans les cas où on laisse $c'$ \deleted{\unsure{et $z$ et $c'$}} fixe, et où on envoie $\underline{\mathrm{Ch}}_{\Sigma,m}\inserted{(X;x,y)}$ dans $\underline{\mathrm{Ch}}_{\Sigma,m+2}(X;x,y)$, ou inversement $c$ fixe et $\underline{\mathrm{Ch}}_{\Sigma,n}(X;y,z)$ dans $\underline{\mathrm{Ch}}_{\Sigma,n+2}$. Prenons le premier cas, et notons un chemin \inserted{$c$} sans indices négatifs (ce n'était pas une notation bien heureuse !)
\diagram{\begin{tikzcd}[ampersand replacement=\&,column sep=.8em]
A_0\arrow[r] \& A_1 \& A_2\arrow[l] \& \cdots \& {}\& A_{2m}\arrow[l]
\end{tikzcd}}
(NB \deleted{\unsure{tous nos}} à cause de la condition $\unsure{n|m}$ pair, tous nos chemins commencent désormais par une flèche de $A_0$ dans $A_1$, jamais en sens inverse) \deleted{Il \unsure{faut} comparer les deux chemins
\diagram{\begin{tikzcd}[ampersand replacement=\&,column sep=.15em,scale=.81,transform shape,cells={font=\tiny}]
A_0\arrow[r,"\mathrm{id}","{\deleted{\unsure{A}}}"'] \& A_0 \& A_0\arrow[l,"\mathrm{id}"']\arrow[r] \& A_1 \& A_2\arrow[l] \& \cdots \& A_{2m}\arrow[l]\arrow[r,"\mathrm{id}"] \& A_{2m}\arrow[r,dash,"\mathrm{id}"] \& \underset{\substack{\Vert\\\unsure{B_0}}}{A_{2m}}\arrow[r,dash,"{\unsure{\to|\leftarrow}}" description] \& B_1 \& B_2\arrow[l] \& \cdots \& B_{2n}\arrow[l]
\end{tikzcd}}}

, Je vois que l'explicitation de l'homotopie, sans offrir de difficulté, est très fastidieuse à expliquer. Je \deleted{préfère} vais essayer de donner une démonstration conceptuelle, sans ``calcul'' fastidieux

% archive sheet number: 15
